%% TEST-NAME: stress-ta-inline-gray
%% TEST-PURPOSE: Inline TA stress fixture with gray backrefs.
%% TEST-KIND: stress
%% TEST-BEHAVIOR: B-NUM-SHARED, B-REND-INLINE, B-LINK-CLICKABLE, B-LINK-NO-ORPHAN
%% TEST-STATUS: EXPLORATORY
%% TEST-RENDER-MODES: inline, gray-backrefs
%% TEST-PDF-PAGES: 3
%% TEST-PDF-STEXT: Stress fixture inline gray
%% TEST-PDF-STEXT: (?s)First theorem for inline joint-atom proof.*\{9\.9, 9\.10\}\*
%% TEST-PDF-STEXT: (?s)Second theorem for inline joint-atom proof.*\{9\.9, 9\.10\}\*

\def\START{}\def\END{}
\documentclass{article}
\usepackage[T1]{fontenc}
\usepackage{amsmath,amssymb}
\usepackage{mathtools}
\usepackage{graphicx}

%% Load the real theorem style from the monograph
\usepackage{sty-theorems-ta}

%% codependent with depth=1 (default, matches numberwithin=section)
\usepackage[backref-style=inline]{codependent}
\codepsetup{backref-color=gray}
\codeptrack{theorem,proposition,lemma,corollary,definition,notation,terminology,example,remark,question}

\usepackage[hypertexnames=false]{hyperref}
\usepackage{cleveref}

%% cleveref names for custom envs that thmtools doesn't auto-register
\crefname{notation}{Notation}{Notations}
\crefname{terminology}{Terminology}{Terminologies}
\crefname{question}{Question}{Questions}

\title{Stress fixture inline gray}
\author{Test}
\date{}

\begin{document}
\START
\maketitle

\section{Counter sharing under load}

Some introductory paragraph to check paragraph numbering.

\begin{definition}\label{def:widget}
A \emph{widget} is a thing that does stuff.
\end{definition}

\begin{theorem}\label{thm:nontrivial}
Every widget (see \cref{def:widget}) is nontrivial.
\end{theorem}

\begin{proof}
By \cref{def:widget}, obvious.
\end{proof}

\begin{corollary}\label{cor:trivial}
By \cref{thm:nontrivial}, no widget is trivial.
\end{corollary}

\begin{remark}\label{rem:note}
This follows from \cref{thm:nontrivial} and \cref{cor:trivial}.
\end{remark}

\begin{proposition}\label{prop:compose}
Widgets compose. This extends \cref{thm:nontrivial}.
\end{proposition}

\begin{lemma}\label{lem:technical}
A technical lemma used in the proof of \cref{prop:compose}.
\end{lemma}

\begin{notation}\label{not:W}
We write $W$ for the category of widgets (\cref{def:widget}).
\end{notation}

\begin{terminology}\label{term:idempotent}
A widget (\cref{def:widget}) that composes with itself
(\cref{prop:compose}) is called \emph{idempotent}.
\end{terminology}

\begin{example}\label{ex:empty}
The empty set is not a widget (\cref{def:widget}).
\end{example}

\begin{question}\label{q:gadget}
Is every widget (\cref{def:widget}) a gadget?
See \cref{thm:nontrivial,cor:trivial,prop:compose}.
\end{question}

A closing paragraph referencing \cref{lem:technical} and \cref{rem:note}.
All ten envs above should be numbered 1.1 through 1.10 sequentially.

\section{Display math suppression}

A paragraph before display math.

\begin{theorem}\label{thm:idempotent}
For all widgets $w$ (\cref{def:widget}), we have
\begin{equation}\label{eq:widget}
  w \circ w = w.
\end{equation}
In particular, $w$ is idempotent (\cref{term:idempotent}).
\end{theorem}

\begin{proof}
By \cref{lem:technical}, consider the aligned computation:
\begin{align}
  w \circ w &= w \circ (w \circ w) \\
            &= w.
\end{align}
\end{proof}

A paragraph after the proof. The equation and align should NOT get
margin atom numbers.

\section{Cross-section backrefs}

\begin{theorem}\label{thm:sec3}
This uses \cref{def:widget} and \cref{thm:nontrivial} from section 1.
Those atoms should show ``Used in 3.1.''
\end{theorem}

\begin{definition}\label{def:gadget}
A \emph{gadget} extends a widget (\cref{def:widget}).
Answers \cref{q:gadget} partially.
\end{definition}

\begin{corollary}\label{cor:sec3}
Combines \cref{thm:sec3}, \cref{def:gadget}, and \cref{prop:compose}.
\end{corollary}

Cross-section remote backref read: see also \codepbackrefsof{thm:nontrivial}
for related results.

A second cross-section remote read: see also \codepbackrefsof{thm:idempotent}
for display-math dependencies.


%% ====================================================================
%%  Section 9: Auto-detect proof scenarios (W05-A2 inline parity)
%%  The inline stress variants are intentionally shorter than the
%%  appendix-gray fixture, so this jumps to section 9 to keep the joint
%%  atom source token aligned with appendix-gray: {9.9, 9.10}*.
%% ====================================================================
\setcounter{section}{8}
\section{Auto-detect proof scenarios}
\label{sec:w05a1}

%% --- 9.1: Backward-ref body-cite ---
\begin{theorem}\label{thm:s9-backward}
Every inline widget (\cref{def:widget}) inherits the technical reduction
of \cref{lem:technical}.
\end{theorem}

We recall the widget definitions before proceeding.

\begin{proof}[Proof of \cref{thm:s9-backward}]
By \cref{def:widget} and the technical lemma \cref{lem:technical}.
\end{proof}

%% --- Forward-ref body-cite: proof comes BEFORE the theorem ---
\begin{proof}[Proof of \cref{thm:s9-forward-late}]
By widget composition (\cref{prop:compose}) and idempotence terminology
(\cref{term:idempotent}).
\end{proof}

\begin{theorem}\label{thm:s9-forward-late}
The forward-referenced inline result combines \cref{prop:compose}
with \cref{term:idempotent}.
\end{theorem}

%% --- 9.3, 9.4: Joint-via-single-cref best-effort serial case ---
\begin{proposition}\label{prop:s9-joint-a}
First inline joint target for best-effort serial proof attribution.
\end{proposition}

\begin{corollary}\label{cor:s9-joint-b}
Second inline joint target for best-effort serial proof attribution.
Follows from \cref{prop:s9-joint-a}.
\end{corollary}

We give a unified proof of both inline results.

\begin{proof}[Proof of \cref{prop:s9-joint-a,cor:s9-joint-b}]
Both \cref{prop:s9-joint-a} and \cref{cor:s9-joint-b} follow from
\cref{lem:technical} and widget composition.
\end{proof}

%% --- 9.5: Joint-via-single-cref mixed-tracking case ---
\begin{lemma}\label{lem:s9-mixed}
Tracked inline target for the mixed-tracking joint proof scenario.
\end{lemma}

The tracked lemma and the section reference play different roles here.

\begin{proof}[Proof of \cref{lem:s9-mixed,sec:w05a1}]
The tracked lemma follows from \cref{def:widget}; the section ref is
navigational and triggers the tracked-entity gate.
\end{proof}

%% --- Section-only heading: zero records expected ---
\begin{proof}[Proof of \autoref{sec:w05a1}]
Section-only heading exercises the tracked-entity gate.  It mentions
\cref{thm:nontrivial} in the body, which is dropped by the gate.
\end{proof}

%% --- 9.9, 9.10: Joint-atom proof through inline Used-In rendering ---
\begin{theorem}\label{thm:s9-jatom-x}
First theorem for inline joint-atom proof.
Uses \cref{def:gadget}.
\end{theorem}

\begin{theorem}\label{thm:s9-jatom-y}
Second theorem for inline joint-atom proof.
Follows from \cref{lem:technical}.
\end{theorem}

We give a joint proof of both inline theorems above.

\begin{proof}[Proof of \cref{thm:s9-jatom-x,thm:s9-jatom-y}]
Both follow from \cref{def:gadget} and \cref{lem:technical} by a unified
inline-mode argument.
\end{proof}

\END
\end{document}
